\begin{abstract}
Test-Time Adaptation (TTA) aims to adapt a source-pretrained model to a shifting target distribution using only unlabeled online data. While the recently proposed COME method advances this setting by incorporating evidential (Dirichlet) uncertainty into a conservative opinion-entropy objective, it treats every test sample uniformly: the same conservative entropy loss is applied regardless of whether the model is genuinely uncertain or already confident. This uniform treatment discards useful supervisory signal that confident samples could otherwise provide and offers no mechanism linking the evidential uncertainty to the softmax confidence used for prediction. We propose \textbf{Uncertainty-Calibrated Dual-Path Adaptation (UCDPA)}, a framework that resolves these limitations through three components. First, we introduce the \emph{Uncertainty-Confidence Balance (UCB) score}, a differentiable routing signal $b_i=\sigma((c_i-\tau_c)/T_c)\,\sigma((1-u_i-\tau_u)/T_u)$ that is monotonically increasing in softmax confidence $c_i$ and monotonically decreasing in Dirichlet uncertainty $u_i$, enabling per-sample soft routing between two complementary objectives. Second, we design a \emph{scale-normalized dual-path} adaptation loss that applies conservative opinion-entropy minimization to uncertain samples and pseudo-label cross-entropy to confident samples, with class-balanced reweighting and a calibration regularizer that maintains batch-mean consistency between confidence and uncertainty. Third, we add a feature prototype alignment term that pulls representations toward their pseudo-label class prototypes. We provide a theorem on UCB monotonicity and a proposition on calibration consistency. On ImageNet-C with ViT-B/16, UCDPA improves mean corruption accuracy by \ucdpadelta{} percentage points over COME (\ucdpaacc{} vs.\ \comeacc{}), with statistically significant gains ($p<\pvalue{}$), lower Expected and Maximum Calibration Error (\ucdpaece{} vs.\ \comeece{}), and improved error-detection AUROC (\ucdpaauroc{} vs.\ \comeauroc{}). We further validate UCDPA on CIFAR-100-C with a CIFAR-ResNet-18 backbone (BatchNorm adaptation), demonstrating cross-backbone and cross-dataset generalization. Ablations confirm that every component contributes, and sensitivity sweeps show the method is stable to routing hyperparameter choices within a \sensRange{}~pp range.
\end{abstract}
